% =====================================================================
%  4.4 Drøfting av funksjoner — teori og oppgaver
% =====================================================================

\del{1}{Fortegnslinje for $f'$ og ekstremalpunkter}

\begin{ordliste}
voksende & $f$ er voksende på et intervall når grafen stiger fra venstre mot høyre. & \mittsprak\\
avtagende & $f$ er avtagende på et intervall når grafen synker fra venstre mot høyre. & \mittsprak\\
ekstremalpunkt & Et topp- eller bunnpunkt. Oppgis som koordinater $(x, f(x))$. & \mittsprak\\
terrassepunkt & Et punkt med $f'(a) = 0$ der $f'$ ikke skifter fortegn. Grafen flater ut og fortsetter samme vei. & \mittsprak\\
drøfte & Finne ut hvordan en funksjon oppfører seg: nullpunkter, voksende og avtagende, ekstremalpunkter. & \mittsprak\\
\end{ordliste}

\begin{mustkunne}
\begin{minipage}[t]{0.50\linewidth}
\textbf{Hvorfor:} Tangenten i $x = a$ har stigningstall $f'(a)$. Er $f'(a) > 0$, stiger tangenten, og grafen stiger. Er $f'(a) < 0$, synker grafen. Fortegnet til $f'$ sier altså hvilken vei $f$ går.\par\smallskip
\textbf{Drøfting i fem steg:}
\begin{enumerate}[leftmargin=1.4em,itemsep=0pt,topsep=2pt]
  \item Deriver og faktoriser $f'(x)$.
  \item Løs $f'(x) = 0$.
  \item Lag fortegnslinje for $f'$ (ikke for $f$).
  \item Sett $x$-verdiene inn i $f$ og oppgi punktene.
  \item Finn nullpunktene (faktoriser, eller finn ett nullpunkt og bruk polynomdivisjon) og skisser grafen.
\end{enumerate}
\end{minipage}\hfill
\begin{minipage}[t]{0.47\linewidth}
\renewcommand{\arraystretch}{1.5}
\begin{tabular}[t]{@{}l l@{}}
$f'(x) > 0$ & $f$ er voksende\\
$f'(x) < 0$ & $f$ er avtagende\\
$f'$ går fra $+$ til $-$ & toppunkt\\
$f'$ går fra $-$ til $+$ & bunnpunkt\\
$f'(a) = 0$, samme fortegn & terrassepunkt\\
\end{tabular}\par\smallskip
Ved et ekstremalpunkt er tangenten vannrett, så $f'(a) = 0$.\par\smallskip
{\footnotesize Figur til eksempel 1: grafen til $f(x) = x^3 - 3x^2$ og fortegnslinjen for $f'$.}\par
\centering
\begin{tikzpicture}[font=\small,fpos/.style={line width=1.1pt},fneg/.style={line width=1.1pt,dash pattern=on 4pt off 3pt}]
\begin{axis}[name=graf, at={(0,0)}, anchor=south west, scale only axis, width=5.4cm, height=3.3cm,
  axis lines=middle, xmin=-1.0769, xmax=3.0769, ymin=-5, ymax=2.5, enlargelimits=false,
  xtick={0,2,3}, ytick={-4}, tick label style={font=\scriptsize}, xlabel={$x$}, ylabel={$y$}, label style={font=\small}]
\addplot[domain=-1.07:3.03, thick, hoved, samples=60] {x^3-3*x^2};
\addplot[only marks, mark=*, mark size=1.6pt, hoved] coordinates {(0,0) (2,-4)};
\node[font=\scriptsize, anchor=south west] at (axis cs:0.1,0.1) {topp};
\node[font=\scriptsize, anchor=north west] at (axis cs:2.1,-4.0) {bunn};
\end{axis}
\begin{scope}[shift={(0,-1.1)}]
  \draw[graa,thin] (1.40,0.40) -- (1.40,-1.95);
  \draw[graa,thin] (4.00,0.40) -- (4.00,-1.95);
  \node[anchor=east] at (-0.15,-0.00) {$x$};
  \draw[fneg] (0.00,-0.00) -- (1.40,-0.00);
  \draw[fpos] (1.40,-0.00) -- (4.00,-0.00);
  \draw[fpos] (4.00,-0.00) -- (5.40,-0.00);
  \node[fill=white,inner sep=1.5pt] at (1.40,-0.00) {$0$};
  \node[anchor=east] at (-0.15,-0.70) {$x - 2$};
  \draw[fneg] (0.00,-0.70) -- (1.40,-0.70);
  \draw[fneg] (1.40,-0.70) -- (4.00,-0.70);
  \draw[fpos] (4.00,-0.70) -- (5.40,-0.70);
  \node[fill=white,inner sep=1.5pt] at (4.00,-0.70) {$0$};
  \node[anchor=east] at (-0.15,-1.55) {$f'(x)$};
  \draw[fpos,hoved] (0.00,-1.55) -- (1.40,-1.55);
  \draw[fneg,hoved] (1.40,-1.55) -- (4.00,-1.55);
  \draw[fpos,hoved] (4.00,-1.55) -- (5.40,-1.55);
  \node[fill=white,inner sep=1.5pt] at (1.40,-1.55) {$0$};
  \node[fill=white,inner sep=1.5pt] at (4.00,-1.55) {$0$};
  \node[below] at (1.40,-1.95) {$0$};
  \node[below] at (4.00,-1.95) {$2$};
\end{scope}
\end{tikzpicture}
\end{minipage}
\end{mustkunne}

\begin{eksempel}[1 — full drøfting]
Drøft $f(x) = x^3 - 3x^2$.

\begin{stegtabell}
$f'(x) = 3x^2 - 6x = 3x(x - 2)$ & Deriver ledd for ledd. Sett $3x$ utenfor.\\
$f'(x) = 0$ for $x = 0$ og $x = 2$ & Nullproduktregelen.\\
Fortegnslinje for $f'$ & Faktorene gir $+, -, +$ (figuren i boksen over). Tallet $3$ er positivt.\\
Voksende: $\langle\leftarrow, 0]$ og $[2, \rightarrow\rangle$ & $f' > 0$ gir voksende.\\
Avtagende: $[0, 2]$ & $f' < 0$ gir avtagende.\\
$f(0) = 0$ og $f(2) = 8 - 12 = -4$ & Sett inn i $f$. Topp $(0, 0)$, bunn $(2, -4)$.\\
$x^2(x - 3) = 0$ gir $x = 0$ og $x = 3$ & Nullpunktene er $(0, 0)$ og $(3, 0)$.\\
\end{stegtabell}
\end{eksempel}

\begin{eksempel}[2 — terrassepunkt]
Drøft $g(x) = x^3 - 6x^2 + 12x$.

\begin{stegtabell}
$g'(x) = 3x^2 - 12x + 12 = 3(x - 2)^2$ & Deriver, sett $3$ utenfor og bruk kvadratsetningen.\\
$g'(x) = 0$ for $x = 2$ & Bare ett nullpunkt for $g'$.\\
$g'(x) = 3(x - 2)^2 > 0$ for alle $x \ne 2$ & Et kvadrat er aldri negativt. Fortegnslinjen viser ingen fortegnsskifte i $x = 2$.\\
$g(2) = 8 - 24 + 24 = 8$ & Terrassepunkt $(2, 8)$. Ikke topp eller bunn.\\
$g$ er voksende for alle $x$ & Grafen flater ut i $(2, 8)$ og stiger videre.\\
$g(x) = x(x^2 - 6x + 12)$ & Andre faktor har $D = -12 < 0$. Eneste nullpunkt er $x = 0$.\\
\end{stegtabell}
\begin{center}
\begin{tikzpicture}[font=\small,fpos/.style={line width=1.1pt},fneg/.style={line width=1.1pt,dash pattern=on 4pt off 3pt}]
  \draw[graa,thin] (1.40,0.40) -- (1.40,-1.95);
  \node[anchor=east] at (-0.15,-0.00) {$x - 2$};
  \draw[fneg] (0.00,-0.00) -- (1.40,-0.00);
  \draw[fpos] (1.40,-0.00) -- (2.80,-0.00);
  \node[fill=white,inner sep=1.5pt] at (1.40,-0.00) {$0$};
  \node[anchor=east] at (-0.15,-0.70) {$x - 2$};
  \draw[fneg] (0.00,-0.70) -- (1.40,-0.70);
  \draw[fpos] (1.40,-0.70) -- (2.80,-0.70);
  \node[fill=white,inner sep=1.5pt] at (1.40,-0.70) {$0$};
  \node[anchor=east] at (-0.15,-1.55) {$g'(x)$};
  \draw[fpos,hoved] (0.00,-1.55) -- (1.40,-1.55);
  \draw[fpos,hoved] (1.40,-1.55) -- (2.80,-1.55);
  \node[fill=white,inner sep=1.5pt] at (1.40,-1.55) {$0$};
  \node[below] at (1.40,-1.95) {$2$};
\end{tikzpicture}
\end{center}
\end{eksempel}

\begin{vanligfeil}
\begin{itemize}[leftmargin=1.2em,itemsep=1pt,topsep=0pt]
  \item «$f'(a) = 0$ gir alltid topp eller bunn» gjelder ikke generelt. For $g(x) = x^3 - 6x^2 + 12x$ er $g'(2) = 0$, men $g'$ er positiv på begge sider, så det er et terrassepunkt.
  \item Fortegnslinje for $f$ i stedet for $f'$ er feil. Fortegnene til $f'(x) = 3x(x-2)$ gir $+, -, +$. Det sier hvilken vei $f$ går, ikke om $f$ er positiv.
  \item «Toppunktet er $x = 0$» er ikke ferdig. Et punkt har to koordinater: $(0, f(0)) = (0, 0)$ for $f(x) = x^3 - 3x^2$.
  \item Å sette $x$-verdien inn i $f'$ for å finne $y$ er feil. Bruk $f$: $y = f(2)$, ikke $f'(2)$.
\end{itemize}
\end{vanligfeil}

\nivaahode{1}

\oppg{1.1}{1}{\uh}
La $f(x) = 2x^3 - 6x$. Fullfør.
\begin{tasks}(1)
  \task $f'(x) = \ \rute\, x^2 - \ \rute$
  \task $f'(x) = 6(x - \ \rute\,)(x + \ \rute\,)$
  \task $f'(x) = 0$ for $x = $ \skrivlinje[2cm] og $x = $ \skrivlinje[2cm]
\end{tasks}

\oppg{1.2}{1}{\uh}
Fortegnslinjen for $f'(x) = (x + 2)(x - 1)$ er tegnet.
\begin{center}
\begin{tikzpicture}[font=\small,fpos/.style={line width=1.1pt},fneg/.style={line width=1.1pt,dash pattern=on 4pt off 3pt}]
  \draw[graa,thin] (1.40,0.40) -- (1.40,-1.95);
  \draw[graa,thin] (3.30,0.40) -- (3.30,-1.95);
  \node[anchor=east] at (-0.15,-0.00) {$x + 2$};
  \draw[fneg] (0.00,-0.00) -- (1.40,-0.00);
  \draw[fpos] (1.40,-0.00) -- (3.30,-0.00);
  \draw[fpos] (3.30,-0.00) -- (4.70,-0.00);
  \node[fill=white,inner sep=1.5pt] at (1.40,-0.00) {$0$};
  \node[anchor=east] at (-0.15,-0.70) {$x - 1$};
  \draw[fneg] (0.00,-0.70) -- (1.40,-0.70);
  \draw[fneg] (1.40,-0.70) -- (3.30,-0.70);
  \draw[fpos] (3.30,-0.70) -- (4.70,-0.70);
  \node[fill=white,inner sep=1.5pt] at (3.30,-0.70) {$0$};
  \node[anchor=east] at (-0.15,-1.55) {$f'(x)$};
  \draw[fpos,hoved] (0.00,-1.55) -- (1.40,-1.55);
  \draw[fneg,hoved] (1.40,-1.55) -- (3.30,-1.55);
  \draw[fpos,hoved] (3.30,-1.55) -- (4.70,-1.55);
  \node[fill=white,inner sep=1.5pt] at (1.40,-1.55) {$0$};
  \node[fill=white,inner sep=1.5pt] at (3.30,-1.55) {$0$};
  \node[below] at (1.40,-1.95) {$-2$};
  \node[below] at (3.30,-1.95) {$1$};
\end{tikzpicture}
\end{center}
\begin{tasks}(1)
  \task Skriv intervallene der $f$ er voksende og der $f$ er avtagende.
  \task For hvilke $x$-verdier har $f$ toppunkt og bunnpunkt? Forklar med fortegnet til $f'$.
\end{tasks}

\oppg{1.3}{1}{\uh}
La $f(x) = x^2 - 6x + 5$. Fullfør.
\begin{tasks}(1)
  \task $f'(x) = \ \rute\, x - \ \rute$, så $f'(x) = 0$ gir $x = $ \skrivlinje[2cm]
  \task $f'$ skifter fra \skrivlinje[1.5cm] til \skrivlinje[1.5cm] i dette punktet. Grafen har et \skrivlinje[2.5cm].
  \task $f(\ \rute\,) = $ \skrivlinje[2cm], så punktet er $(\ \rute\,,\ \rute\,)$
\end{tasks}

\nivaahode{2}

\oppg{1.4}{2}{\uh}
La $f(x) = x^3 + 3x^2 - 9x$.
\begin{tasks}(1)
  \task Finn $f'(x)$ og faktoriser. Lag fortegnslinje for $f'$.
  \task Skriv intervallene der $f$ er voksende og avtagende.
  \task Finn ekstremalpunktene og oppgi dem som koordinater.
\end{tasks}

\oppg{1.5}{2}{\uh}
La $g(x) = x^3 - 3x^2 - 9x + 11$.
\begin{tasks}(1)
  \task Vis at $x = 1$ er et nullpunkt. Del $g(x)$ på $x - 1$ med polynomdivisjon. Finn alle nullpunktene eksakt.
  \task Finn ekstremalpunktene med fortegnslinje for $g'$.
  \task Skisser grafen. Marker nullpunkter, ekstremalpunkter og skjæring med $y$-aksen.
\end{tasks}

\oppg{1.6}{2}{\uh}
La $f(x) = x^4 - 4x^3$. Finn alle $x$ der $f'(x) = 0$. Avgjør hva slags punkt hvert av dem er, og oppgi punktene som koordinater.

\oppg{1.7}{2}{\uh}
Finn \textbf{den første feilen}. Løs riktig.
\[ \begin{aligned}
f(x) = x^3 - 3x^2 - 24x \ &\longrightarrow\ f'(x) = 3x^2 - 6x - 24 \ \longrightarrow\ f'(x) = 3(x - 4)(x + 2)\\
&\longrightarrow\ x = 4 \text{ eller } x = -2\\
&\longrightarrow\ f' \text{ er negativ for } x < -2,\ \text{positiv for } -2 < x < 4,\ \text{negativ for } x > 4\\
&\longrightarrow\ \text{bunnpunkt } (-2, 28),\ \text{toppunkt } (4, -80)
\end{aligned} \]

\nivaahode{3}

\oppg{1.8}{3}{\uh}
Er påstanden alltid, noen ganger eller aldri sann? Begrunn svaret.
\begin{tasks}(1)
  \task Hvis $f$ er voksende på et intervall, er $f'(x) > 0$ for alle $x$ i intervallet.
  \task Et andregradspolynom har nøyaktig ett ekstremalpunkt.
  \task Hvis $f$ og $g$ har samme fortegnslinje for den deriverte, er $f(x) = g(x)$.
\end{tasks}

\oppg{1.9}{3}{\uh}
Lag selv et polynom $f$ med toppunkt i $x = -1$ og bunnpunkt i $x = 2$.
\begin{tasks}(1)
  \task Skriv $f$ og vis at det virker.
  \task Hvor mange slike polynomer finnes? Forklar.
\end{tasks}
\hint{For hvilke $x$ må $f'$ være null?}

\begin{sjekk}
\begin{tasks}(2)
  \task Finn bunnpunktet til $f(x) = x^2 - 8x + 3$.
  \task $f'$ skifter fra $+$ til $-$ i $x = 3$. Hva slags punkt har $f$ der?
  \task! Hvorfor viser $f'(x) > 0$ at grafen til $f$ stiger? Forklar med tangenter.
\end{tasks}
\end{sjekk}

% ---------------------------------------------------------------------

\del{2}{Største og minste verdi, optimering}

\begin{ordliste}
største verdi & Det høyeste $y$-tallet grafen har på intervallet. & \mittsprak\\
minste verdi & Det laveste $y$-tallet grafen har på intervallet. & \mittsprak\\
endepunkt & Intervallets ytterpunkter. På $[a, b]$ er $a$ og $b$ endepunkter. & \mittsprak\\
optimering & Å finne den beste verdien, for eksempel størst volum eller minst kostnad. & \mittsprak\\
\end{ordliste}

\begin{mustkunne}
\begin{minipage}[t]{0.57\linewidth}
\textbf{På et lukket intervall $[a, b]$} ligger største og minste verdi enten der $f'(x) = 0$ eller i et endepunkt. Hvorfor: Der $f' \ne 0$ stiger eller synker grafen, så det finnes både en høyere og en lavere verdi like ved. Bare i vannrette punkter og ved endene stopper grafen.\par\smallskip
\textbf{Metode:} Regn ut $f$ i endepunktene og i alle $x$ med $f'(x) = 0$ inni intervallet. Sammenlign tallene. Er endepunktene ikke med, bruker du fortegnslinja for $f'$.\par\smallskip
\textbf{Optimering i fire steg}
\begin{enumerate}[leftmargin=1.4em,itemsep=0pt,topsep=2pt]
  \item Velg variabel og finn definisjonsmengden.
  \item Sett opp funksjonen.
  \item Deriver og finn kandidatene.
  \item Sammenlign og svar i en setning med enhet.
\end{enumerate}
\end{minipage}\hfill
\begin{minipage}[t]{0.40\linewidth}
\vspace{0pt}\centering
\begin{tikzpicture}[baseline=(current bounding box.north)]
\begin{axis}[width=\linewidth, height=1.15\linewidth, axis lines=middle, xmin=-0.5, xmax=5.8, ymin=-3, ymax=24,
  xtick={0,5}, ytick={4,20}, tick label style={font=\scriptsize}, xlabel={$x$}, ylabel={$y$}, label style={font=\small}]
\addplot[domain=0:5, thick, hoved, samples=60] {x^3-6*x^2+9*x};
\addplot[only marks, mark=*, mark size=1.6pt, hoved] coordinates {(0,0) (1,4) (3,0) (5,20)};
\node[font=\scriptsize, anchor=south] at (axis cs:1,5.5) {lokal topp};
\node[font=\scriptsize, anchor=south east] at (axis cs:5.0,21.5) {største verdi};
\node[font=\scriptsize, anchor=north] at (axis cs:3,-1.2) {bunn};
\end{axis}
\end{tikzpicture}\par
{\footnotesize Største verdi på $[0, 5]$ ligger i endepunktet.}
\end{minipage}
\end{mustkunne}

\begin{eksempel}[1 — største og minste verdi]
Finn største og minste verdi til $f(x) = x^3 - 3x^2$ på $[-1, 4]$.

\begin{stegtabell}
$f'(x) = 3x(x - 2) = 0$ for $x = 0$ og $x = 2$ & Begge ligger inni $[-1, 4]$. Det er kandidatene.\\
$f(-1) = -1 - 3 = -4$ & Endepunkt.\\
$f(0) = 0$ & Toppunkt.\\
$f(2) = 8 - 12 = -4$ & Bunnpunkt.\\
$f(4) = 64 - 48 = 16$ & Endepunkt.\\
Største verdi: $16$ (i $x = 4$) & Sammenlign de fire tallene. Største verdi ligger i et endepunkt, ikke i toppunktet.\\
Minste verdi: $-4$ (i $x = -1$ og $x = 2$) & Samme verdi i endepunktet og bunnpunktet. Minste verdi kan nås to steder.\\
\end{stegtabell}
\end{eksempel}

\begin{eksempel}[2 — optimering: esken fra 3.3]
Nå kan vi løse det vi lovte i 3.3. En eske lages av et papp på $20$ cm $\times$ $30$ cm. Vi klipper bort kvadrater med side $x$ cm i hjørnene og bretter opp sidene. Finn $x$ som gir størst volum. Bruk hjelpemidler.

\begin{stegtabell}
1. $0 \le x \le 10$ & Variabel $x$ cm. Kortsiden gir $20 - 2x \ge 0$. Definisjonsmengden er $[0, 10]$.\\
2. $V(x) = x(20 - 2x)(30 - 2x)$ & Lengde $\cdot$ bredde $\cdot$ høyde.\\
$V(x) = 4x^3 - 100x^2 + 600x$ & Gang ut.\\
3. $V'(x) = 12x^2 - 200x + 600 = 0$ & Deriver. Del på $4$: $3x^2 - 50x + 150 = 0$.\\
$x = \dfrac{50 \pm \sqrt{700}}{6}$, så $x \approx 3,92$ eller $x \approx 12,74$ & Andregradsformelen. Verdien $12,74$ ligger utenfor $[0, 10]$ og forkastes.\\
4. $V(0) = 0$, $V(10) = 0$, $V(3,92) \approx 1056$ & Vi tar med endepunktene (flat eske, $V = 0$) for å kunne sammenligne.\\
\end{stegtabell}
\textbf{Svar:} Esken får størst volum, ca. $1056$ cm$^3$, når kvadratene i hjørnene har side ca. $3,92$ cm.
\end{eksempel}

\begin{vanligfeil}
\begin{itemize}[leftmargin=1.2em,itemsep=1pt,topsep=0pt]
  \item «Største verdi er alltid i et toppunkt» gjelder ikke generelt. Et toppunkt er bare lokalt størst. Grafen i boksen over har toppunkt $(1, 4)$, men største verdi $20$ på $[0, 5]$.
  \item Å glemme endepunktene er feil. For $f(x) = x^2$ på $[1, 3]$ gir $f'(x) = 0$ bare $x = 0$, utenfor intervallet. Største verdi $f(3) = 9$ finnes bare ved å sjekke endepunktet.
  \item Å bruke en løsning av $f'(x) = 0$ utenfor intervallet er feil. I eskeoppgaven er $x \approx 12,74$ umulig, siden pappet bare er $20$ cm bredt.
  \item Å svare «$x = 3,92$» på «finn størst volum» er ikke nok. Svar på det som spørres om: volum med enhet.
\end{itemize}
\end{vanligfeil}

\nivaahode{1}

\oppg{2.1}{1}{\uh}
Finn største og minste verdi til $f(x) = x^2 - 4x$ på $[0, 5]$. Fullfør.
\begin{tasks}(1)
  \task $f'(x) = \ \rute\, x - \ \rute$, så $f'(x) = 0$ for $x = $ \skrivlinje[2cm]
  \task $f(0) = $ \skrivlinje[1.5cm] \quad $f(\ \rute\,) = $ \skrivlinje[1.5cm] \quad $f(5) = $ \skrivlinje[1.5cm]
  \task Største verdi er \skrivlinje[1.5cm] og minste verdi er \skrivlinje[1.5cm].
\end{tasks}

\oppg{2.2}{1}{\uh}
Finn største og minste verdi.
\begin{tasks}(1)
  \task $f(x) = x^3 - 12x$ på $[-3, 3]$
  \task $g(x) = -x^2 + 4x$ på $[3, 6]$
\end{tasks}

\oppg{2.3}{1}{\uh}
Du har $40$ m gjerde og skal lage en rektangulær innhegning mot en husvegg. Det trengs gjerde på tre sider. De to sidene vinkelrett på veggen er $x$ m. Fullfør de fire stegene.
\begin{tasks}(1)
  \task Siden langs veggen er $40 - 2x$ m. Definisjonsmengde: $0 \le x \le$ \skrivlinje[1.5cm]
  \task $A(x) = x(40 - 2x) = $ \skrivlinje[4cm]
  \task $A'(x) = $ \skrivlinje[3cm] $= 0$ gir $x = $ \skrivlinje[1.5cm]
  \task Regn ut $A(0)$, $A(\ \rute\,)$ og $A(20)$. Svar i en setning med enhet.
\end{tasks}

\nivaahode{2}

\oppg{2.4}{2}{\uh}
En åpen eske uten lokk har kvadratisk bunn med side $x$ cm og høyde $h$ cm. Hele overflaten er $300$ cm$^2$. Bruk de fire stegene til å finne største volum.
\begin{tasks}(1)
  \task Vis at $h = \dfrac{300 - x^2}{4x}$ og at $V(x) = 75x - \dfrac{x^3}{4}$. Finn definisjonsmengden.
  \task Finn $V'(x)$ og løs $V'(x) = 0$.
  \task Vis med fortegnslinje for $V'$ at $V(10)$ er den største verdien. Svar i en setning med enhet.
\end{tasks}

\oppg{2.5}{2}{\uh}
Finn \textbf{den første feilen}. Løs riktig.
Største og minste verdi til $f(x) = 2x^3 - 9x^2 + 12x$ på $[0, 4]$:
\[ \begin{aligned}
f'(x) = 6x^2 - 18x + 12 \ &\longrightarrow\ f'(x) = 6(x - 1)(x - 2) \ \longrightarrow\ x = 1 \text{ eller } x = 2\\
&\longrightarrow\ f(1) = 5 \text{ og } f(2) = 4\\
&\longrightarrow\ \text{største verdi } 5,\ \text{minste verdi } 4
\end{aligned} \]

\oppg{2.6}{2}{\uh}
Konsentrasjonen av et legemiddel i blodet er $K(t) = -0,5t^3 + 4,5t^2$ mg/L etter $t$ timer, for $0 \le t \le 8$ (tenkt modell).
\begin{tasks}(1)
  \task Finn $K'(t)$. Når er konsentrasjonen størst, og hva er den da? Svar med enhet.
  \task Finn den laveste konsentrasjonen for $1 \le t \le 8$. Hvor ligger den, og hvorfor ikke i et toppunkt eller bunnpunkt?
  \task I hvilket tidsrom avtar konsentrasjonen?
\end{tasks}

\nivaahode{3}

\oppg{2.7}{3}{\uh}
\begin{tasks}(1)
  \task For hvilke $b$ er den største verdien til $f(x) = x^3 - 3x$ på $[-2, b]$ lik $2$? Begrunn svaret.
  \task Er påstanden alltid, noen ganger eller aldri sann? Begrunn svaret. «Hvis $f$ er voksende på $[a, b]$, er $f(b)$ den største verdien til $f$ på $[a, b]$.»
\end{tasks}

\oppg{2.8}{3}{\uh}
La $f(x) = x^3 - 27x + k$ på $[0, 6]$, der $k$ er et tall.
\begin{tasks}(1)
  \task Forklar hvorfor $x$-verdiene der $f$ har størst og minst verdi ikke avhenger av $k$.
  \task Finn $k$ slik at største verdi er $100$. Hva er minste verdi da?
\end{tasks}

\begin{sjekk}
\begin{tasks}(2)
  \task Finn største og minste verdi til $f(x) = x^2 - 6x$ på $[1, 5]$.
  \task Et rektangel har sidene $x$ og $12 - x$. Skriv arealfunksjonen og definisjonsmengden.
  \task! Hvorfor må vi regne ut $f$ i endepunktene, selv om $f'(x) \ne 0$ der?
\end{tasks}
\end{sjekk}
